\section{Introduction}
\label{sec:introduction}

Approximate homomorphic encryption (FHE) makes encrypted real-number computation practical by allowing controlled approximation error.
That error is also relevant in the context of security analysis.
Li and Micciancio showed that observing suitably chosen approximate decryptions can reveal the secret key of CKKS-like schemes~\cite{lm,ckks}; related attacks now reach aggressively parameterized exact FHE as well~\cite{cheon2024attacks}.
In approximate FHE one cannot simply require every decrypted result to equal an exact message: approximation is the advertised interface.

Li, Micciancio, Schultz-Wu, and Sorrell (LMSS) formalized this attack vector through the \emph{IND-CPAD} security game, and analyzed a noise-flooding defense~\cite{lmss}.
The defense perturbs each admissible decryption answer with a fresh, sufficiently wide discrete Gaussian.
On the other hand, the decryption result is useful only when that Gaussian width is small enough to preserve precision, so this parameter choice has itself become a research topic \cite{costache2023precision, bergamaschi2025revisiting}.

The parameter trade-off also exposes a mechanization problem.
Suppose a one-query replacement costs \(\delta\) in statistical distance, then a standard game hop argument (i.e. using data processing inequality) would pay \(q\delta\).
The LMSS analysis instead bounds a conditional KL divergence at each query and applies the Micciancio--Walter Pythagorean preservation argument only after constructing the whole transcript~\cite{mw17}, paying \(O(\sqrt q)\) in statistical distance instead.
For fixed target advantage, this is the difference between choosing flooding width proportional to \(\sqrt q\) and to \(q\), with a corresponding impact on precision and FHE parameters.

We verify that security argument in Rocq on top of SSProve, a framework for machine-checked, game-based cryptography~\cite{ssprove}.
The top-level theorem has ordinary cryptographic semantics: it bounds the winning probability of a concrete IND-CPAD game by the IND-CPA bound of a constructed reduction plus an explicit loss.
Along the way, we construct a new relational program logic and verify its soundness, and implement a verified compiler that transforms arbitrary SSProve code into those compatible with our program logic.

\subsection{Result and security significance}

Let \(S\) be an approximate-FHE scheme whose message space has local charts in
\(\bbZ^n\), let \(\mathsf{NF}_{\gamma}[S]\) be its noise-flooded transform, and
let \(\cA\) make at most \(q\) non-failing decryption queries in the IND-CPAD
game.
Under the explicit assumptions summarized in
\Cref{tab:theorem-boundary}, the formal development constructs an IND-CPA
adversary \(\cB_{\cA,q}\) and proves
\begin{equation}
\label{eq:headline}
 \WinProb[\mathsf{IND\text{-}CPAD}_{\mathsf{NF}_{\gamma}[S]}(\cA)]
 \leq
 \beta_{\mathsf{CPA}}(\cB_{\cA,q})+
 \frac{\sqrt{qn}}{2\gamma}.
\end{equation}
The expression \(\beta_{\mathsf{CPA}}\) is the supplied winning-probability
bound for the underlying scheme.
The formal theorem is given in the artifact as
\texttt{NoiseFloodingSecure.is\_secure}.

The result covers adaptive oracle programs rather than adversaries prewritten as \(q\) separate query phases.
Its proof isolates one lossy transition: real flooded decryption is replaced by flooding around the plaintext recorded in the challenge table.
Every other change---opening and relinking packages, compiling calls, replacing unreachable residual decryptions, and identifying the reduction game---is proved exact.

\subsection{Why a program logic?}

SSProve's original architecture connects two proof styles: high-level equational reasoning about cryptographic packages and low-level relational reasoning about their procedures~\cite{ssprove}.
Noise flooding needs a new quantitative connection of the same kind.
The local fact is a KL-divergence bound between two implementations of one decryption call; the security claim is an additive bound between two games containing an adaptive adversary.

An ordinary approximate-equivalence judgment does not preserve the required parameter dependence.
Its sequence rule adds scalar errors, so applying Pinsker at every oracle call turns a one-call cost \(\sqrt{\epsilon/2}\) into \(q\sqrt{\epsilon/2}\).
Our Pythagorean judgment instead records a tuple of conditional KL costs.
Sequencing concatenates tuples, and a checked Micciancio--Walter rule converts the accumulated tuple to additive error only once.
After \(q\) calls of KL cost \(\epsilon\), this gives \(\sqrt{q\epsilon/2}\).

The remaining obstacle is adaptivity: a concrete SSProve adversary is one stateful oracle program, not a syntactic list of \(q\) calls.
We therefore verify a compiler that exposes selected oracle calls and prove a generic local-to-adaptive lifting theorem.
Together, the judgment and compiler form a reusable program-logic interface: prove one invariant-preserving oracle rule, then obtain a whole-program additive bound for arbitrary well-typed adaptive code.
Noise flooding is the first substantial application of this interface; \Cref{thm:compile-replace} states the general theorem.

\subsection{Contributions}

We separate the cryptographic result from the machinery used to establish it.

\emph{A checked noise-flooding reduction.} We define the IND-CPA and IND-CPAD games, the noise-flooded construction, and a concrete reduction as SSProve packages.
Rocq checks the complete hybrid chain and the loss in \Cref{eq:headline}.
The paper presents the game argument in cryptographic terms; selected formal definitions appear in the appendices.

\emph{A relational program logic for SSProve.}
We define \emph{additive-error} and \emph{Pythagorean} judgments, and prove
several inference rules in Rocq.
The logic combines unary invariant reasoning,
direct coupling and additive-error rules, KL sampling rules, tuple-concatenating
sequencing, and a final bridge to statistical distance.
% Missing
% subdistribution mass is an explicit outcome, so assertions cannot make a
% comparison spuriously easier.

\emph{A verified local-to-adaptive oracle rule.} We verify a compiler that exposes selected calls in an arbitrary concrete SSProve program.
Its main theorem lifts one invariant-preserving Pythagorean oracle judgment to a whole-program additive-error judgment with loss \(\sqrt{q\norm{s}_1/2}\).
This is a program-logic result independent of the particular FHE game.

\emph{Probability and discrete-Gaussian analysis.} We verify KL divergence and its finiteness obligations, Pinsker's inequality, a conditional-coordinate preservation theorem, and the equal-variance KL formula for integer discrete Gaussians.
These results provide the semantic soundness argument and the derived vector-Gaussian sampling rule used by the logic.

\emph{An auditable theorem boundary.} The 27,986-line artifact is organized so that games, assumptions, loss, and final theorem can be reviewed separately from proof scripts.
We report its trusted base, distinguish inherited library assumptions from local proof obligations, and provide a verified replacement for the one unverified real analysis lemma inherited through the current SSProve stack.

\subsection{Theorem boundary and conventional obligations}
\label{sec:scope}

The formally verified theorem roughly says that, given any \emph{approximately correct} and \emph{IND-CPA secure} FHE scheme, noise flooding can transform such a scheme to become IND-CPAD secure.
In particular, we assume the underlying FHE scheme includes deterministic decryption, approximate correctness with certainty (i.e. no bad key pairs exist), and a particular metric space structure for plaintexts.
Our formulation cleanly isolates one step in the entire FHE security analysis.
For an imperfectly correct scheme such as CKKS, one must add a standard game hop to capture the corresponding correctness error; this does not alter our Pythagorean hop.
We have not yet formalized this additional ``up-to-bad" game hop.

For the current stage of the project, we also assume the message space is locally isometric to $\bbZ^n$.
This captures the usual instantiation (i.e. polynomial rings over ideals) while being simpler to audit.

We have not yet formally verified any concrete FHE scheme to be IND-CPA secure or approximately correct.
Instead, we discuss the CKKS scheme in \Cref{sec:limitations}, and leave its formalization and proof engineering for future work.
As a result, we do not claim a verified deployment or concrete parameter set.

Our result also follows SSProve's usual convention: probabilistic polynomial time is not formalized in the Rocq/SSProve semantics~\cite{ssprove}, and should be checked by the usual external audit.
For the LMSS noise flooding, the standard reduction runs \(\cA\) once and adds oracle forwarding, table bookkeeping, and Gaussian sampling; no unusual complexity issue is hidden here.

\subsection{Road map}

\Cref{sec:security-model} defines the construction, attack game, and theorem.
\Cref{sec:verified-reduction} presents the game hops.
\Cref{sec:pyth-toolkit} presents the program logic and its local-to-adaptive lifting theorem, and \Cref{sec:artifact} gives the audit map and trusted base.
